\documentclass[10pt,letterpaper]{article}
\usepackage[
    lmargin=2.2in,
    rmargin=0.7in,
    tmargin=1.0in,
    bmargin=1.0in
]{geometry}
\usepackage{notes}

\begin{document}

\begin{center}
\Large\textbf{DAY 30 -- Eigenvalues, Eigenvectors, and Linear Systems}\\
\normalsize (Section 4.3)
\end{center}

%%%%%%%%%%%%%%%%%%%%%%%%%%%%%%%%%%%%%%%%%%%%%%%%%%%%%%%%%%%%%%%%%%%%%%%%
\section{Finding Solutions to a Linear System}

\subsection{Motivation}

Our goal for the lecture is to solve systems of differential equations
of the form
\begin{equation} \label{1}
\mathbf{x}'=A\mathbf{x}, \quad \textbf{x}(0)=\textbf{x}_0
\end{equation}
where $A$ is a constant $2\times2$ matrix and $\mathbf{x}(t)$ is a vector of unknown functions.


%-------------------------------------------------------------
\subsection{An Exponential Guess}

Solving the matrix problem $\mathbf{x}'=A\mathbf{x}$ has many similarities to solving $x'=4x$.

\cols{
\begin{center}
    \textbf{1D Case}
\end{center}
\[
x'=4x
\]
We want a function whose derivative is a constant
multiple of itself. An exponential is a natural guess:
\[
x(t)=e^{\lambda t}.
\]
Substituting gives
\begin{align*}
x'&=4x,\\
\lambda e^{\lambda t}&=4e^{\lambda t}
\end{align*}
Thus, $\lambda=4$. The general solution is
\[
x(t)=Ce^{4t}.
\]
}{
\begin{center}
    \textbf{2D Case}
\end{center}
\[
\mathbf{x}'=A\mathbf{x}
\]
We guess an exponential multiplied by a constant vector:
\[
\mathbf{x}(t)=e^{\lambda t}\mathbf{v},
\qquad \mathbf{v}\ne\mathbf{0}.
\]
Substituting gives
\begin{align*}
\mathbf{x}'&=A\mathbf{x},\\
\lambda e^{\lambda t}\mathbf{v}
&=e^{\lambda t}A\mathbf{v},\\
\lambda\mathbf{v}&=A\mathbf{v}.
\end{align*}
If we find such a scalar and nonzero vector, we obtain solutions
$$
    \textbf{x}_1(t) = C e^{\lambda t} \textbf{v}. 
$$
}

Solving this equation 
\begin{equation}    
    A\mathbf{v} = \lambda \mathbf{v}, \quad \mathbf{v}\ne\mathbf{0}
\end{equation}
is called the \textbf{eigenvalue problem}. We must learn how to solve the eigenvalue problem before we can solve the system \eqref{1}.


%%%%%%%%%%%%%%%%%%%%%%%%%%%%%%%%%%%%%%%%%%%%%%%%%%%%%%%%%%%%%%%%%%%%%%%%
\section{Review of Fundamental Concepts}

% Separate the algebra review from the exponential motivation.
\subsection{Matrix Operations and Determinants}
Consider the matrices
\[
A=\begin{pmatrix}1&2\\3&4\end{pmatrix},
\qquad I=\begin{pmatrix}1&0\\0&1\end{pmatrix}.
\]
Let $\lambda$ be any number. We can define $M:=A-\lambda I$ and calculate its components
\[
M:=A-\lambda I
=\begin{pmatrix}1&2\\3&4\end{pmatrix}
-\begin{pmatrix}\lambda&0\\0&\lambda\end{pmatrix}
=\begin{pmatrix}1-\lambda&2\\3&4-\lambda\end{pmatrix}.
\]
Only the diagonal entries change. Next, recall the definition of a determinant 
$$
\text{If } A=\begin{pmatrix}
    a & b\\
    c & d
\end{pmatrix} \text{ then } \det(A) = \begin{vmatrix}
    a & b\\
    c & d
\end{vmatrix}= ad-bc.
$$
We can calculate the determinant of the $M$ matrix 
\begin{align*}
\det(M)=\det(A-\lambda I)=& \begin{vmatrix}
    1-\lambda & 2\\
    3 & 4-\lambda
\end{vmatrix}\\
&=(1-\lambda)(4-\lambda)-2(3)\\
&=\lambda^2-5\lambda-2.
\end{align*}
The result is a quadratic polynomial in $\lambda$.

\subsection{Nonhomogeneous Systems}
An algebraic system $$M\mathbf{v}=\mathbf{b}$$ is
\textbf{nonhomogeneous} when $\mathbf{b}\ne\mathbf{0}$.
If $\det(M)\ne0$, it has exactly one solution. For example,
\[
\begin{pmatrix}1&2\\3&4\end{pmatrix}
\begin{pmatrix}u\\w\end{pmatrix}
=\begin{pmatrix}3\\7\end{pmatrix}
\quad\Longleftrightarrow\quad
\begin{dcases}u+2w=3,\\3u+4w=7.\end{dcases}
\]
The determinant is $4-6=-2\ne0$, so there is exactly one solution. The solution is $u=1$, $w=1.$ There are no other combinations $(u,w)$ that satisfy the system.

\subsection{Homogeneous Systems}
A system is \textbf{homogeneous} when its right-hand side is zero:
\[
M\mathbf{v}=\mathbf{0}.
\]
The zero vector $\textbf{v}=\textbf{0}$ always solves this system; hence, we call it the
\textbf{trivial solution}. Any nonzero solution is called
\textbf{nontrivial}.
\begin{itemize}
\item If $\det(M)\ne0$, the only solution is $\mathbf{v}=\mathbf{0}$.
\item If $\det(M)=0$, there are infinitely many solutions,
including nontrivial solutions.
\end{itemize}
A homogeneous system can never have no solution. If you were asked to find one answer out of infinite possible choices, it is wise to choose the simplest one.

\begin{example}[0]{    
Find one non-trivial solution to 
\[
\begin{pmatrix}1&2\\3&6\end{pmatrix}
\begin{pmatrix}x\\y\end{pmatrix}=\begin{pmatrix}0\\0\end{pmatrix}.
\]
}
First, notice that coefficient matrix has a zero determinant. The second equation is three times the first, so they impose the same condition:
\[
x+2y=0,\qquad 3x+6y=0.
\]
Any point on the line $y=-\tfrac{1}{2}x$ is a valid answer. We can choose any value of $x$, except $x=0$. We can choose $x=2$, for example, then we get
\[
\begin{pmatrix}
    x\\y
\end{pmatrix}=\begin{pmatrix}
 x\\
 -\tfrac{1}{2}x
\end{pmatrix}
=
\begin{pmatrix}2\\-1\end{pmatrix}.
\]
That's it. $(2,-1)$ is a nontrivial answer.
\end{example}

\subsection{Connecting to the Eigenvalue Problem}
To solve the eigenvalue problem, we first convert it to a \textit{homogeneous} system:
\[
A\mathbf{v}=\lambda\mathbf{v}
\quad\Longleftrightarrow\quad
(A-\lambda I)\mathbf{v}=\mathbf{0}.
\]
We need $\mathbf{v}\ne\mathbf{0}$. That means we cannot allow $\det(A-\lambda I)\ne0$ because that only has the trivial solution. Therefore, in order for a nontrivial $\textbf{v}$ to exist, we require 
\[
\boxed{\det(A-\lambda I)=0.}
\]
This is the \textbf{characteristic equation}. Its roots are the
\textbf{eigenvalues}. For each eigenvalue, we then find a nonzero
solution of $(A-\lambda I)\mathbf{v}=\mathbf{0}$, called an
\textbf{eigenvector}. Next, we discuss how to find the eigenvalues and eigenvectors.

%%%%%%%%%%%%%%%%%%%%%%%%%%%%%%%%%%%%%%%%%%%%%%%%%%%%%%%%%%%%%%%%%%%%%%%%
\section{Eigenvalues and Eigenvectors}


\subsection{Finding Eigenvalues}
In order to find the eigenvalues $\lambda$ for a matrix $A$, you must first set up the characteristic equation
$$
\det(A-\lambda I)=0
$$
and then solve for $\lambda$. Since this is a quadratic equation, we expect \textit{two values}. Let's call them
$$
\lambda_1  \quad \text{and} \quad \lambda_2.
$$
Here is a demonstration.

\begin{example}[1]{
Find the two eigenvalues of
\[
A=
\begin{pmatrix}
4&1\\
2&3
\end{pmatrix}.
\]
}

\paragraph{Step 1. Form the matrix $A-\lambda I$.}
Subtract $\lambda$ from each diagonal entry:
\[
M=A-\lambda I
=
\begin{pmatrix}
4&1\\2&3
\end{pmatrix}
-
\begin{pmatrix}
\lambda&0\\0&\lambda
\end{pmatrix}
=
\begin{pmatrix}
4-\lambda&1\\
2&3-\lambda
\end{pmatrix}.
\]

\paragraph{Step 2. Compute the determinant.} Calculate $\det(M)$. It is standard to rearrange the terms in the form $a\lambda^2 + b\lambda + c$
\begin{align*}
\det(A-\lambda I)
&=(4-\lambda)(3-\lambda)-2\\
&=12-3\lambda -4\lambda +\lambda^2-2\\
&=\lambda^2-7\lambda+10.
\end{align*}

\paragraph{Step 3. Solve the characteristic equation.} Sometimes, we need the quadratic formula to solve $a\lambda^2 + b\lambda + c=0$. But here, it is easier to factor the polynomial
\[
\lambda^2-7\lambda+10
=(\lambda-5)(\lambda-2)=0.
\]
Therefore,
\[
\boxed{\lambda_1=5,\qquad \lambda_2=2.}
\]
\end{example}

%-------------------------------------------------------------
\subsection{An Equivalent Form}

If you want to find eigenvalues, you may alternatively use
\[
\boxed{\det(\lambda I-A)=0} \quad  \text{instead of} \quad \det(A-\lambda I)=0
\]
This gives the same characteristic polynomial. Here is a redo of Example 1:
\begin{align*}
\det(\lambda I-A)
&=
\begin{vmatrix}
\lambda-4&-1\\
-2&\lambda-3
\end{vmatrix}\\
&=(\lambda^2-7\lambda+12)-2\\
&=\lambda^2-7\lambda+10.
\end{align*}
I personally prefer using $\lambda I-A$ over $A-\lambda I$ because it puts $\lambda$ first
in the diagonal entries. It is easier to mentally expand $(\lambda-4)(\lambda-3)$ rather than $(4-\lambda)(3-\lambda)$.

%-------------------------------------------------------------
\begin{example}[2]{
Find the eigenvalues of
\[
A=
\begin{pmatrix}
0&1\\
-2&-3
\end{pmatrix}.
\]
}

\paragraph{Step 1. Form $\lambda I-A$.}
\[
\lambda I-A=
\begin{pmatrix}
\lambda&-1\\
2&\lambda+3
\end{pmatrix}.
\]

\paragraph{Step 2. Compute the determinant.}
\[
\det(\lambda I-A)
=\lambda(\lambda+3)+2
=\lambda^2+3\lambda+2.
\]

\paragraph{Step 3. Solve the characteristic equation.}
\[
\lambda^2+3\lambda+2
=(\lambda+1)(\lambda+2)=0.
\]
Therefore,
\[
\boxed{\lambda_1=-1,\qquad \lambda_2=-2.}
\]
\end{example}


%-------------------------------------------------------------

\subsection{Finding Eigenvectors}
Once an eigenvalue is known, substitute it into
$(A-\lambda I)\mathbf{v}=\mathbf{0}$ and choose one nonzero solution. Since there are infinite choices for eigenvectors, simply choose the one with the cleanest numbers.

\begin{example}[3]{
For the matrix from Example 1,
\[
A=
\begin{pmatrix}
4&1\\
2&3
\end{pmatrix},
\]
find an eigenvector corresponding to each eigenvalue
$\lambda_1=5$ and $\lambda_2=2$.
}

\paragraph{1. Substitute $\lambda_1=5$ into $A-\lambda I$.}
Substituting $\lambda = 5$ into $\begin{pmatrix}
4-\lambda&1\\
2&3-\lambda
\end{pmatrix}$ gives
\[
A-5I=
\begin{pmatrix}
    4-5& 1 \\
    2 & 3-5
\end{pmatrix}=
\begin{pmatrix}
-1&1\\
2&-2
\end{pmatrix}.
\]
\paragraph{2. Set up and solve $(A-5 I)\textbf{v}=\textbf{0}$.}
We can let $\mathbf{v}=\begin{pmatrix}u\\w\end{pmatrix}$. Then
the equation $(A-5I)\mathbf{v}=\mathbf{0}$ becomes
\[
\begin{dcases}
-u+w=0,\\
2u-2w=0.
\end{dcases}
\]
Both equations say $w=u$. There are infinitely many answers. The easiest is to choose $u=1$, which means $w=1$. So the eigenvector for $\lambda_1=5$ is
\[
\boxed{
\mathbf{v}_1=\begin{pmatrix}
    u\\w
\end{pmatrix}=\begin{pmatrix}1\\1\end{pmatrix}.
}
\]

\paragraph{3. Substitute $\lambda_2=2$ into $A-\lambda I$.}
Substitution yields 
\[
A-2I=\begin{pmatrix}
    4-2& 1 \\
    2 & 3-2
\end{pmatrix}=
\begin{pmatrix}
2&1\\
2&1
\end{pmatrix}.
\]
\paragraph{4. Set up and solve $(A-2 I)\textbf{v}=\textbf{0}$.}
The equation $(A-2I)\mathbf{v}=\mathbf{0}$ becomes
\[
\begin{dcases}
2u+w=0,\\
2u+w=0.
\end{dcases}
\]
Thus, $w=-2u$. To keep it simple, choose $u=1$, which means $w=-2$. Therefore, the eigenvector for $\lambda_2=2$ is 
\[
\boxed{
\mathbf{v}_2=\begin{pmatrix}1\\-2\end{pmatrix}.
}
\]
In these examples, we can choose either component and then use the equation to determine the other. The two components cannot be chosen independently, and the resulting vector must be nonzero.
\end{example}

%-------------------------------------------------------------
\subsection{Properties of Eigenvectors}

Keep three facts in mind:
\begin{enumerate}
    \item \textbf{Nonzero}: an eigenvector must be nontrivial: $\mathbf{v}\ne\mathbf{0}$.

    \item \textbf{Correspondence}: an eigenvector belongs
    to a particular eigenvalue. Keep each eigenvalue
    paired with its eigenvector.

    \item \textbf{Nonuniqueness}: if $\mathbf{v}$ is an
    eigenvector, then $c\mathbf{v}$ is also an eigenvector
    for the same eigenvalue whenever $c\ne0$.
\end{enumerate}
Individual components may be zero; the entire vector
cannot be zero.

%-------------------------------------------------------------
\subsection{Tip on Choosing Simple Components}

Choose components that avoid unnecessary fractions.
For example, suppose you want to find the cleanest solution to 
\[
3u+7w=0.
\]
Choosing $u=1$ works, but it is not clever because then $w=-3/7$. In contrast, 
choosing $u=7$ gives $w=-3$. Thus, we obtain a convenient 
eigenvector with integers
\[
\mathbf{v}=\begin{pmatrix}7\\-3\end{pmatrix}.
\]
Here is a tip: suppose you want to find the eigenvector satisfying 
\[
\begin{pmatrix}
    a & b\\
    a & b
\end{pmatrix} \textbf{v} = \begin{pmatrix}
    0\\0
\end{pmatrix}
\]
You may choose your eigenvector to be 
\[
\boxed{\mathbf{v}=\begin{pmatrix}b\\-a\end{pmatrix}.}
\]

\subsection{The General Theory}

The goal of the eigenvalue problem is to find a number $\lambda$ and a
nonzero vector $\mathbf{v}$ such that
\[
A\mathbf{v}=\lambda\mathbf{v}.
\]
There will be two values  of $\lambda$ which are called the \textbf{eigenvalue} of $A$. Each $\lambda$ has a vector $\mathbf{v}\ne\mathbf{0}$ is its \textbf{eigenvector}. Together, $(\lambda,\mathbf{v})$ form an \textbf{eigenpair}.
We exclude the zero vector because $A\mathbf{0}=\lambda\mathbf{0}$ holds for every
$\lambda$ and tells us nothing about $A$. To find the two values of $\lambda$, you must solve the quadratic equation
$$
    \det(A-\lambda I)=0 \quad \text{or} \quad \det(\lambda I-A)=0.
$$
Then you must find an eigenvector for each eigenvalue by finding one nontrivial solution to 
$$
    (A-\lambda I)\textbf{v}=\textbf{0} \quad \text{or} \quad (\lambda I-A) \textbf{v}=0.
$$

%%%%%%%%%%%%%%%%%%%%%%%%%%%%%%%%%%%%%%%%%%%%%%%%%%%%%%%%%%%%%%%%%%%%%%%%

\section{Solving the System}

\subsection{The Initial Value Problem}

Finally, we are ready to solve systems of DEs. Suppose you are given a linear system of DEs
\begin{equation} \label{2}    
\textbf{x}' = \begin{pmatrix}
    a & b\\
    c & d
\end{pmatrix} \textbf{x}
\end{equation}
and an initial condition 
$$
\textbf{x}(0) = \begin{pmatrix}
    x_0\\y_0
\end{pmatrix}. 
$$
Your task is to find a parametric solution $\textbf{x}(t)=(x(t),y(t))$ that satisfies both the system and the IC. It comes down to memorizing a formula.

\subsection{General Solution Formula}

Here is the formula for the general solution to \eqref{2}:
\[
\boxed{
\mathbf{x}(t)
=c_1e^{\lambda_1t}\mathbf{v}_1
+c_2e^{\lambda_2t}\mathbf{v}_2.
}
\]
In this formula, $\lambda_1, \textbf{v}_1$ and $\lambda_2, \textbf{v}_2$ are the two eigenpairs of the matrix $A$. The constant coefficients $c_1$ and $c_2$ are arbitrary. Notice how similar this formula is to the solution
$$
    x(t) = c_1 e^{\lambda_1 t}+c_2 e^{\lambda_2 t}.
$$
The values of $c_1$ and $c_2$ can be determined by the IC.  Let's first focus on finding the general solution.

%-------------------------------------------------------------
\begin{example}[4]{
Find the general solution to
\[
\mathbf{x}'=
\begin{pmatrix}
4&1\\
2&3
\end{pmatrix}\mathbf{x}.
\]
}

\paragraph{Step 1. Find both eigenpairs.}
We have already found the two eigenpairs for the matrix $A$. From Examples 1 and 3, we have
\[
\lambda_1=5,\qquad
\mathbf{v}_1=\begin{pmatrix}1\\1\end{pmatrix},
\]
and
\[
\lambda_2=2,\qquad
\mathbf{v}_2=\begin{pmatrix}1\\-2\end{pmatrix}.
\]

\paragraph{Step 2. Write the general solution.} Next, you write down the formula which you have memorized, and then substitute the corresponding values.
\begin{align*}    
\mathbf{x}(t)&=c_1e^{\lambda_1t}\mathbf{v}_1 +c_2e^{\lambda_2t}\mathbf{v}_2\\
&=
\boxed{
c_1e^{5t}\begin{pmatrix}1\\1\end{pmatrix}
+c_2e^{2t}\begin{pmatrix}1\\-2\end{pmatrix}.
}
\end{align*}
Not everyone's general solution may look the same due to the non-uniqueness of eigenvectors. For example, it is valid to choose $\textbf{v}_1=(2,2)$ instead of $\textbf{v}_1=(1,1)$ as the eigenvector. However, the arbitrary constant $c_1$ negates any constant multiple differences between eigenvectors.

\end{example}

%-------------------------------------------------------------
\subsection{Applying Initial Conditions}

Suppose we are given
\[
\mathbf{x}(0)=\mathbf{x}_0.
\]
Since both exponential factors equal $1$ at $t=0$,
\[
c_1\mathbf{v}_1+c_2\mathbf{v}_2=\mathbf{x}_0.
\]
The vectors $\textbf{v}_1$ and $\textbf{v}_2$ are known. If
\[
\mathbf{v}_1=\begin{pmatrix}u_1\\w_1\end{pmatrix},
\qquad
\mathbf{v}_2=\begin{pmatrix}u_2\\w_2\end{pmatrix},
\]
then solving for the coefficients amounts to solving a linear system 
\[
\begin{pmatrix}
u_1&u_2\\
w_1&w_2
\end{pmatrix}
\begin{pmatrix}c_1\\c_2\end{pmatrix}
=\mathbf{x}_0.
\]
To solve an IVP, you have to solve a system at the end. This is the reason why you want to choose eigenvectors with clean numbers.

%-------------------------------------------------------------
\begin{example}[5]{
Solve the IVP:
\[
\mathbf{x}'=
\begin{pmatrix}
4&1\\
2&3
\end{pmatrix}\mathbf{x},
\qquad
\mathbf{x}(0)=\begin{pmatrix}7\\1\end{pmatrix}.
\]
}

\paragraph{Step 1. Write the general solution.}
From Example 4,
\[
\mathbf{x}(t)
=c_1e^{5t}\begin{pmatrix}1\\1\end{pmatrix}
+c_2e^{2t}\begin{pmatrix}1\\-2\end{pmatrix}.
\]

\paragraph{Step 2. Apply the initial condition.}
At $t=0$,
\[
c_1\begin{pmatrix}1\\1\end{pmatrix}
+c_2\begin{pmatrix}1\\-2\end{pmatrix}
=\begin{pmatrix}7\\1\end{pmatrix}.
\]
Equating components gives
\[
\begin{dcases}
c_1+c_2=7,\\
c_1-2c_2=1.
\end{dcases}
\]
Subtracting the second equation from the first,
\[
3c_2=6
\implies c_2=2.
\]
Then $c_1=7-2=5$.

\paragraph{Step 3. State the solution.}
\[
\boxed{
\mathbf{x}(t)
=5e^{5t}\begin{pmatrix}1\\1\end{pmatrix}
+2e^{2t}\begin{pmatrix}1\\-2\end{pmatrix}.
}
\]
\end{example}

%-------------------------------------------------------------
\subsection{Putting the Steps Together}
If you want to solve an IVP $\textbf{x}'=A\textbf{x}$ with $\textbf{x}(0)=\textbf{x}_0$, then you must find eigenpairs and solve for the unknown coefficients. Here is a full example from beginning to end.

\begin{example}[6]{
Solve the IVP:
\[
\mathbf{x}'=
\begin{pmatrix}
3&1\\
0&2
\end{pmatrix}\mathbf{x},
\qquad
\mathbf{x}(0)=\begin{pmatrix}4\\2\end{pmatrix}.
\]
}

\paragraph{Step 1. Find the eigenvalues.}
The characteristic equation is
\[
\det(\lambda I-A)
=
\begin{vmatrix}
\lambda-3&-1\\
0&\lambda-2
\end{vmatrix}
=(\lambda-3)(\lambda-2)=0.
\]
Therefore,
\[
\lambda_1=3,\qquad \lambda_2=2.
\]

\paragraph{Step 2. Find the eigenvectors.}
For $\lambda_1=3$,
\[
(A-3I)\mathbf{v}=\mathbf{0}
\quad\Longleftrightarrow\quad
\begin{pmatrix}
0&1\\
0&-1
\end{pmatrix}
\begin{pmatrix}u\\w\end{pmatrix}
=\begin{pmatrix}0\\0\end{pmatrix}.
\]
We don't have freedom over $w$. Both equations require $w=0$. But, $u$ is completely free. For simplicity, we can choose $u=1$, so the eigenvector for $\lambda_1=3$ is
\[
\mathbf{v}_1=\begin{pmatrix}1\\0\end{pmatrix}.
\]

For $\lambda_2=2$,
\[
(A-2I)\mathbf{v}=\mathbf{0}
\quad\Longleftrightarrow\quad
\begin{pmatrix}
1&1\\
0&0
\end{pmatrix}
\begin{pmatrix}u\\w\end{pmatrix}
=\begin{pmatrix}0\\0\end{pmatrix}.
\]
Thus, $u+w=0$. We have freedom over $u$ or $w$. If we force $u=1$, then $w$ must be $-1$. The 2nd eigenvector is
\[
\mathbf{v}_2=\begin{pmatrix}1\\-1\end{pmatrix}.
\]

\paragraph{Step 3. Write the general solution.} Just memorize the formula and substitute
\[
\mathbf{x}(t)=c_1e^{\lambda_1t}\mathbf{v}_1 +c_2e^{\lambda_2t}\mathbf{v}_2
=c_1e^{3t}\begin{pmatrix}1\\0\end{pmatrix}
+c_2e^{2t}\begin{pmatrix}1\\-1\end{pmatrix}.
\]

\paragraph{Step 4. Apply the initial condition.}
At $t=0$,
\[
c_1\begin{pmatrix}1\\0\end{pmatrix}
+c_2\begin{pmatrix}1\\-1\end{pmatrix}
=\begin{pmatrix}4\\2\end{pmatrix}.
\]
Equating components,
\[
\begin{dcases}
c_1+c_2=4,\\
-c_2=2.
\end{dcases}
\]
Therefore, $c_2=-2$ and $c_1=6$. The solution is
\[
\boxed{
\mathbf{x}(t)
=6e^{3t}\begin{pmatrix}1\\0\end{pmatrix}
-2e^{2t}\begin{pmatrix}1\\-1\end{pmatrix}.
}
\]
\end{example}

\textbf{Note:} This formula 
$$
\mathbf{x}(t)
=c_1e^{\lambda_1t}\mathbf{v}_1
+c_2e^{\lambda_2t}\mathbf{v}_2.
$$
holds for when $A$ has 2 \textit{distinct, real} eigenvalues. There are two other formulas you will need for when the eigenvalues are repeated or when they are complex.

\end{document}